% ============================================================================
% COURS 14 — LOGIQUE, SED ET COMMUNICATION
% Partie B : systèmes à événements discrets et diagrammes SysML
% Source : 18-Logique_SED.docx
% ============================================================================

\section{Systèmes à événements discrets}
\label{sec:sed}

\subsection{Du combinatoire au séquentiel}

Considérons l'ouverture motorisée d'un hayon. Les entrées sont la commande $t_0$ et le capteur de position haute $p_h$. La sortie $M^+$ commande le moteur d'ouverture. Le hayon doit continuer à s'ouvrir après le relâchement de la commande jusqu'à atteindre la position haute.

Une fonction purement combinatoire ne suffit donc pas : il faut mémoriser l'état présent de la sortie. L'équation d'évolution peut s'écrire :
\[
{M^+}_{n+1}=\overline{p_h}\left(t_0+{M^+}_n\right).
\]

\TLLogicImage[.72\linewidth]{image55.jpeg}
\TLLogicImage[.40\linewidth]{image57.jpeg}

\begin{TLLogicDefinition}{Système séquentiel}
Dans un système séquentiel, les sorties dépendent des entrées présentes, de l'état interne antérieur et éventuellement du temps. Une même combinaison d'entrées peut donc conduire à des sorties différentes selon l'histoire du système.
\end{TLLogicDefinition}

La table de vérité est remplacée par une \textbf{table d'états}. Elle relie l'état futur $Q_{n+1}$ à l'état présent $Q_n$ et aux entrées :
\[
Q_{n+1}=f(Q_n,e_1,\ldots,e_p).
\]

Les systèmes à événements discrets se distinguent des systèmes continus : leur évolution est constituée de changements d'état ponctuels déclenchés par des événements.

\subsection{Diagramme d'état SysML}

Le diagramme d'état, noté \texttt{stm}, décrit le comportement d'un bloc SysML. Il représente les états stables du système et les transitions qui les relient.

\begin{center}\TLGrapheEtatsSimple\end{center}

Un diagramme d'état doit être interprété de manière fonctionnelle : il décrit ce que fait le système et dans quelles conditions il change de comportement, sans imposer nécessairement la technologie de réalisation.

\TLLogicImage[.74\linewidth]{image61.png}

\begin{TLLogicMethod}{Construire un graphe d'états}
\begin{enumerate}
  \item Définir la frontière du système et les variables observées ou commandées.
  \item Recenser les situations stables et les nommer par des états.
  \item Tracer les transitions possibles entre états.
  \item Associer à chaque transition un événement, une garde et éventuellement un effet.
  \item Associer aux états les activités d'entrée, continues ou de sortie.
  \item Ajouter l'état initial, les éventuels états finaux et les pseudo-états nécessaires.
  \item Vérifier la complétude, l'absence de contradiction et les scénarios limites.
\end{enumerate}
\end{TLLogicMethod}

\subsection{État et activités associées}

Un état est représenté par un rectangle aux coins arrondis. Il peut contenir un nom et des activités :
\begin{itemize}
  \item \texttt{entry} : activité exécutée une fois lors de l'entrée dans l'état ;
  \item \texttt{do} : activité exécutée tant que l'état reste actif, éventuellement interruptible ;
  \item \texttt{exit} : activité exécutée lors de la sortie de l'état.
\end{itemize}

\begin{center}\TLGrapheEtatDetaille\end{center}

Un état vide représente une attente. Si aucun mot réservé n'est précisé devant une activité, elle est généralement interprétée comme une activité de type \texttt{do}.

\begin{TLLogicApplication}{Lave-linge}
Un cycle simple comporte les états prélavage, lavage, rinçage, essorage et arrêt.
\begin{enumerate}
  \item Identifier l'état initial et les transitions principales.
  \item Associer à chaque état les actionneurs actifs.
  \item Ajouter les temporisations et les conditions de niveau d'eau.
  \item Prévoir une transition prioritaire en cas d'ouverture de porte ou de défaut.
\end{enumerate}
\TLLogicImage[.30\linewidth]{image63.jpeg}
\end{TLLogicApplication}

\subsection{Transitions}

Une transition est instantanée. Elle n'est évaluée que si son état source est actif. Sa notation générale est :
\[
\boxed{\text{événement }[\text{garde}]\,/\,\text{effet}.}
\]

L'état de départ est l'état source et l'état d'arrivée l'état cible. Une transition réflexive quitte l'état puis y revient : les activités \texttt{exit} et \texttt{entry} sont donc exécutées.

\subsubsection{Événement}

Un événement correspond à une occurrence ponctuelle : appui sur un bouton, front d'un capteur, dépassement d'un seuil, réception d'un message. Il n'est pas mémorisé s'il n'est pas traité au moment où il survient.

Les événements temporels usuels sont :
\begin{itemize}
  \item \texttt{after(T)} : déclenchement après une durée $T$ passée dans l'état ;
  \item \texttt{at(D)} : déclenchement à une date absolue $D$ ;
  \item \texttt{when(condition)} : déclenchement lorsque la condition interne change et devient vraie.
\end{itemize}

Sans événement explicite, la fin de l'activité \texttt{do} peut constituer l'événement implicite.

\subsubsection{Garde}

La garde est une condition logique évaluée lorsque l'événement se produit. Contrairement à l'événement, elle décrit un état qui persiste dans le temps. Sans garde, la condition est considérée vraie.

Exemple :
\[
\text{appuiDépart }[\text{porteFermée}\cdot\overline{défaut}].
\]

Une garde doit être mutuellement exclusive avec les autres gardes issues du même choix, sauf lorsqu'une priorité est explicitement définie.

\subsubsection{Effet}

Un effet est une action instantanée réalisée au franchissement de la transition, par exemple :
\[
/\,N:=0,
\qquad
/\,alarme:=1.
\]
Il se distingue des activités associées aux états, qui durent pendant l'activation d'un état.

\begin{TLLogicApplication}{Feu tricolore}
Construire le graphe d'états d'un feu tricolore comportant les états rouge, vert et orange. Ajouter :
\begin{itemize}
  \item les temporisations nominales ;
  \item un mode clignotant orange ;
  \item un retour sécurisé au rouge après un défaut ;
  \item les activités de commande des lampes dans chaque état.
\end{itemize}
\TLLogicImage[.16\linewidth]{image65.jpeg}
\end{TLLogicApplication}

\subsection{Pseudo-états}

Un pseudo-état ne possède pas d'activité. Il organise les chemins du graphe :
\begin{itemize}
  \item initial : point d'entrée obligatoire dans une région ;
  \item final : fin de l'exécution d'une région ;
  \item choix : sélection dynamique d'un chemin selon des gardes ;
  \item jonction : factorisation ou réunion de chemins ;
  \item fork : séparation vers plusieurs régions parallèles ;
  \item join : synchronisation de plusieurs régions ;
  \item historique : reprise du dernier sous-état actif d'un état composite.
\end{itemize}

\begin{center}\TLPseudoEtats\end{center}

\begin{TLLogicWarning}{Choix incomplet ou contradictoire}
Pour un pseudo-état de choix, les gardes sortantes doivent couvrir tous les cas possibles et ne pas être vraies simultanément, sauf priorité explicitement définie. Une branche \texttt{else} permet de garantir la complétude.
\end{TLLogicWarning}

\subsection{États composites}

Un état composite contient des sous-états. Dans un état composite séquentiel, un seul sous-état de la région est actif à la fois.

\begin{center}\TLEtatComposite\end{center}

Une transition entrant directement dans l'état composite active son pseudo-état initial, sauf si elle cible explicitement un sous-état ou un pseudo-état historique.

Le pseudo-état historique superficiel mémorise le dernier sous-état actif au premier niveau. L'historique profond mémorise l'ensemble de la configuration imbriquée.

\subsection{États composites orthogonaux}

Un état composite orthogonal est découpé en plusieurs régions actives simultanément. Chaque région possède sa propre machine à états.

\begin{center}\TLEtatOrthogonal\end{center}

La sortie de l'état composite peut nécessiter une synchronisation par un pseudo-état de jonction. Cette représentation convient aux comportements parallèles : commande et surveillance, déplacement et signalisation, régulation et diagnostic.

\begin{TLLogicMethod}{Vérifier un graphe d'états}
\begin{enumerate}
  \item Chaque région possède un pseudo-état initial.
  \item Aucun état ne doit être inaccessible.
  \item Chaque situation doit conduire à un comportement défini.
  \item Les gardes concurrentes doivent être non contradictoires.
  \item Les boucles sans temporisation ou sans événement doivent être évitées.
  \item Les priorités de sécurité doivent être explicites.
  \item Les régions parallèles doivent être synchronisées avant une sortie commune.
\end{enumerate}
\end{TLLogicMethod}

\begin{TLLogicApplication}{Démarrage d'une turbine à gaz}
Le moteur de lancement entraîne la turbine jusqu'à la vitesse d'autonomie. Le cycle comprend l'embrayage, l'accélération, l'allumage, la détection de la vitesse autonome et le débrayage.
\begin{enumerate}
  \item Identifier les états du cycle et leurs activités.
  \item Compléter les événements et gardes à partir des capteurs de vitesse et d'embrayage.
  \item Construire le chronogramme des variables $E$, $D$, $AV$ et $C$.
  \item Déterminer le temps total de démarrage et le comparer au cahier des charges.
\end{enumerate}
\TLLogicImage[.70\linewidth]{image101.png}
\end{TLLogicApplication}

\section{Diagramme de séquence}
\label{sec:sequence}

Le diagramme de séquence, noté \texttt{sd}, décrit dans l'ordre chronologique les messages échangés entre acteurs, blocs ou composants pour un scénario précis.

\begin{center}\TLDiagrammeSequenceSimple\end{center}

Les éléments essentiels sont :
\begin{itemize}
  \item la ligne de vie de chaque participant ;
  \item la période d'activité, représentée par une bande verticale ;
  \item le message synchrone ou asynchrone ;
  \item le message de retour, souvent en pointillés ;
  \item le message réflexif ;
  \item les fragments combinés.
\end{itemize}

\subsection{Messages}

Un message est une interaction dirigée d'un émetteur vers un récepteur. Son arrivée peut constituer un événement pour le diagramme d'état du récepteur. Un message synchrone bloque l'émetteur jusqu'au retour ; un message asynchrone permet la poursuite de l'exécution.

\subsection{Fragments combinés}

Les fragments encadrent une portion du scénario :
\begin{itemize}
  \item \texttt{alt} : alternatives conditionnelles ;
  \item \texttt{opt} : séquence optionnelle ;
  \item \texttt{loop} : répétition tant qu'une garde est vraie ;
  \item \texttt{par} : séquences exécutées en parallèle ;
  \item \texttt{break} : interruption du reste du scénario.
\end{itemize}

\begin{center}\TLFragmentsSequence\end{center}

\begin{TLLogicApplication}{Vente de repas en ligne}
Décrire le scénario nominal entre client, interface web, service de paiement et cuisine : sélection du repas, validation du panier, paiement, confirmation et préparation.
\begin{enumerate}
  \item Placer les lignes de vie et les messages.
  \item Ajouter un fragment \texttt{alt} pour le paiement accepté ou refusé.
  \item Ajouter un fragment \texttt{loop} pour la sélection de plusieurs articles.
  \item Ajouter un message asynchrone de notification lorsque le repas est prêt.
\end{enumerate}
\TLLogicImage[.64\linewidth]{image108.png}
\end{TLLogicApplication}
